\section{Related Work}\label{sec:related}

Our framework builds upon three research streams that, to date, have
progressed largely independently. We review each in turn, identifying the
specific gaps that motivate our unified approach.

\subsection{Deep Learning and MIL for Lung Adenocarcinoma Subtyping}

Multiple instance learning with attention-based pooling
\cite{ilse2018attention,campanella2019clinical}, coupled with histology-specific
foundation models such as Phikon \cite{filiot2023scaling}, now constitutes the
standard pipeline for weakly-supervised whole-slide image classification.
Within lung adenocarcinoma specifically, several recent efforts have targeted
fine-grained histological subtyping. HybridNet \cite{du2024hybridnet} combines
a CNN and a Transformer to classify five growth patterns including MIP,
achieving 94\% per-subtype accuracy. Miura et al.\ \cite{miura2024establishment}
extended subtyping to seven categories with spatial distribution analysis,
while Huang et al.\ \cite{huang2025attention} employed a Swin Transformer with
multi-scale token learning for the same task.

Despite steadily improving accuracy, these methods share a structural
limitation: they are designed and evaluated exclusively for aggregate metrics
such as overall accuracy and AUC. Under the severe class imbalance inherent
to MIP---where the aggressive subtype occupies only a small fraction of each
slide---optimizing for average performance systematically favors the majority
negative class. Critically, none of these pipelines provides a statistical
guarantee on the false negative rate (FNR), nor a principled mechanism to
identify and defer cases for which the model's prediction is unreliable.
Our Experiments (Section~4.3) further demonstrate that alternative MIL
aggregators---CLAM-SB, DS-MIL, DTFD-MIL, and TransMIL---yield nearly
identical FNR (all around 16\%) to our baseline, confirming that
architectural variation alone cannot resolve this issue.

Our framework retains the state-of-the-art MIL backbone for representation
quality, but augments it with explicit FNR control and selective abstention,
shifting the evaluation criterion from average accuracy to worst-case safety
on the high-risk positive class.


\subsection{Risk-Aware Learning: From Cost-Sensitive to Neyman--Pearson}

Moving beyond average-case optimization requires training and calibration
strategies that explicitly encode the clinical asymmetry of errors.
Cost-sensitive learning (CSL) \cite{el2010foundations,he2009learning} is the
most direct approach: assigning higher penalties to false negatives shifts
the decision boundary so that the model becomes more sensitive to the
positive class, thereby reducing missed detections. However, CSL relies on
heuristically chosen class weights and provides no formal guarantee on the
resulting error rates---a practitioner cannot determine, before deployment,
whether a chosen weight will meet a target FNR on unseen data.

The Neyman--Pearson (NP) paradigm offers a rigorous alternative. NP
classification \cite{rigollet2011neyman} minimizes the type II error (FNR)
subject to a high-probability constraint on the type I error. The NP
umbrella algorithm \cite{tong2018neyman} converts any scoring classifier
into an NP classifier via order-statistic thresholding on a held-out
calibration set of positive samples, yielding a distribution-free,
finite-sample FNR guarantee. Subsequent work has extended NP to multi-class
settings \cite{tian2025neyman}, established minimax optimal rates
\cite{kalan2024distribution}, and explored robustness under label noise
\cite{yao2023asymmetric}. In a parallel direction, risk-calibrated learning
\cite{riskcalibrated2026} embeds clinical severity into the training loss to
suppress critical false negatives, though without the post-hoc guarantee
that NP provides.

Despite these advances, NP calibration has not been integrated into a
MIL-based computational pathology pipeline, nor combined with selective
classification to handle cases where even a calibrated model cannot
confidently decide. Our Phase~I--II design addresses this gap: CSL reshapes
the decision boundary during training to increase sensitivity to the
positive class, and NP calibration subsequently provides a verifiable,
user-tunable FNR bound at inference time.


\subsection{Selective Classification and Conformal Prediction in Medicine}

Even with a calibrated FNR guarantee, some cases remain fundamentally
ambiguous---for example, slides with borderline micropapillary morphology
or imaging artifacts. A clinically safe system must recognize such cases
and defer them to a human pathologist.

Selective classification \cite{hendrickx2024survey} enables a model to
abstain on uncertain inputs, trading coverage for higher reliability on
accepted cases. SelectiveNet \cite{geifman2019selectivenet} jointly trains
a prediction head and a selection head under a coverage constraint, learning
to abstain during training rather than relying on post-hoc thresholding of
softmax scores---which is known to be poorly calibrated under class
imbalance \cite{hendrycks2016baseline}. Conformal prediction
\cite{angelopoulos2022conformal} provides a complementary, post-hoc
framework for calibrating rejection thresholds on held-out data with
distribution-free coverage guarantees.

Several recent works have applied these techniques to medical imaging.
Olsson et al.\ \cite{olsson2022estimating} demonstrated that conformal
prediction reduces diagnostic errors in prostate pathology by an order of
magnitude and detects distribution shifts across scanners. Mehrtens et al.\
\cite{mehrtens2023benchmarking,mehrtens2025pitfalls} benchmarked multiple
uncertainty estimation methods on histopathological WSIs, finding that
selective rejection reliably improves accuracy on both in-distribution and
shifted data, though they also cautioned that conformal methods can be
unreliable under strong covariate shift. For lung cancer, TRUECAM
\cite{zhang2024truecam} integrates conformal prediction with
out-of-distribution detection for NSCLC subtyping from WSIs, and Fayjie
et al.\ \cite{fayjie2024predictive} evaluated uncertainty-based referral
under realistic domain shifts in lung adenocarcinoma classification.

A critical distinction separates these works from ours. All existing
selective classification and conformal methods in medical imaging minimize
a symmetric risk measure on the accepted subset, weighting false positives
and false negatives equally. In MIP detection, the clinical priority is
starkly asymmetric: missing a positive case carries far greater consequence
than raising a false alarm. Our Case~2 and Case~4 formulations
(Section~3.4) explicitly select the abstention threshold to bound the
FNR---rather than the overall error rate---on the accepted subset, while
Case~4 further incorporates cost-sensitive pretraining to recover coverage
lost to aggressive FNR targets. To our knowledge, this is the first work
that combines a SelectiveNet-style architecture with FNR-centric threshold
calibration and cost-sensitive pretraining in a unified pipeline, yielding
a selective classifier that simultaneously controls missed diagnoses,
maximizes automated throughput, and defers ambiguous cases to human
pathologists.


\begin{table}[htbp]
    \centering
    \caption{Conceptual comparison of our work with representative methods
    from each research stream. A $\checkmark$ indicates that the method
    addresses the corresponding requirement; ``$\sim$'' denotes empirical
    mitigation without formal statistical guarantee; ``$^\dagger$'' indicates
    a conformal coverage guarantee that does not specifically target the
    false negative rate.}
    \label{tab:related_comparison}
    \resizebox{\linewidth}{!}{
    \begin{tabular}{lcccc}
        \toprule
        \textbf{Method / Study} & \textbf{MIP Subtyping} & \textbf{FNR Guarantee} & \textbf{Selective Abstention} & \textbf{Statistical Guarantee} \\
        \midrule
        Du et al.\ (HybridNet, 2024)  & \checkmark & -- & -- & -- \\
        Miura et al.\ (2024)          & \checkmark & -- & -- & -- \\
        Huang et al.\ (Swin, 2025)    & \checkmark & -- & -- & -- \\
        \midrule
        CSL \cite{el2010foundations}      & -- & $\sim$ & -- & -- \\
        NP Umbrella \cite{tong2018neyman} & -- & \checkmark & -- & \checkmark \\
        Risk-Calibrated \cite{riskcalibrated2026} & -- & $\sim$ & -- & -- \\
        \midrule
        SelectiveNet \cite{geifman2019selectivenet} & -- & -- & \checkmark & -- \\
        TRUECAM \cite{zhang2024truecam}   & \checkmark & -- & \checkmark & \checkmark$^\dagger$ \\
        Olsson et al.\ \cite{olsson2022estimating} & -- & -- & \checkmark & \checkmark$^\dagger$ \\
        \midrule
        \textbf{This work}           & \checkmark & \checkmark & \checkmark & \checkmark \\
        \bottomrule
    \end{tabular}
    }
\end{table}

As Table~\ref{tab:related_comparison} summarizes, existing work addresses at
most two of the four requirements for clinically safe MIP detection. Our
framework is the only approach that simultaneously satisfies all four.


% ============================================================
% New BibTeX entries required for the rewritten Related Work
% Add these to references.bib before compiling
% ============================================================

% --- MIP Subtyping (Section 2.1) ---

@article{huang2025attention,
  title     = {Attention-driven multi-scale analysis for accurate tumor classification in digital pathology},
  author    = {Huang, Junze and Carello, Saverio J. and others},
  journal   = {Journal of Clinical Oncology},
  year      = {2025},
  volume    = {43},
  number    = {16\_suppl},
  pages     = {e13703},
  doi       = {10.1200/jco.2025.43.16\_suppl.e13703},
  note      = {ASCO 2025 Annual Meeting Abstract},
}

% --- Risk-Aware Learning (Section 2.2) ---

@misc{kalan2024distribution,
  title        = {Distribution-Free Rates in Neyman-Pearson Classification},
  author       = {Mohammadreza M. Kalan and Samory Kpotufe},
  year         = {2024},
  eprint       = {2402.09560},
  archiveprefix = {arXiv},
  primaryclass = {cs.LG},
  doi          = {10.48550/arXiv.2402.09560},
}

@article{yao2023asymmetric,
  author    = {Shunan Yao and Bradley Rava and Xin Tong and Gareth James},
  title     = {Asymmetric Error Control Under Imperfect Supervision: A Label-Noise-Adjusted {Neyman-Pearson} Umbrella Algorithm},
  journal   = {Journal of the American Statistical Association},
  volume    = {118},
  number    = {543},
  pages     = {1824--1836},
  year      = {2023},
  doi       = {10.1080/01621459.2021.2016423},
}

@inproceedings{riskcalibrated2026,
  title     = {Risk-Calibrated Learning: Minimizing Fatal Errors in Medical AI},
  author    = {Abolfazl Mohammadi-Seif and Rajni Rajan},
  booktitle = {Proceedings of the 2026 International Joint Conference on Neural Networks (IJCNN)},
  year      = {2026},
  note      = {arXiv:2604.12693},
}

% --- Selective Classification and Conformal Prediction (Section 2.3) ---

@article{olsson2022estimating,
  author    = {Olsson, Henrik and Kartasalo, Kimmo and Mulliqi, Nita and Capuccini, Marco and Ruusuvuori, Pekka and Samaratunga, Hemamali and Delahunt, Brett and Lindskog, Cecilia and Janssen, Emiel A. M. and Blilie, Anders and {ISUP Prostate Imagebase Expert Panel} and Egevad, Lars and Spjuth, Ola and Eklund, Martin},
  title     = {Estimating diagnostic uncertainty in artificial intelligence assisted pathology using conformal prediction},
  journal   = {Nature Communications},
  year      = {2022},
  volume    = {13},
  number    = {1},
  pages     = {7761},
  doi       = {10.1038/s41467-022-34945-8},
}

@article{mehrtens2023benchmarking,
  author    = {Mehrtens, Hendrik A. and Kurz, Alexander and Bucher, Tabea-Clara and Brinker, Titus J.},
  title     = {Benchmarking common uncertainty estimation methods with histopathological images under domain shift and label noise},
  journal   = {Medical Image Analysis},
  volume    = {89},
  pages     = {102914},
  year      = {2023},
  doi       = {10.1016/j.media.2023.102914},
}

@misc{mehrtens2025pitfalls,
  title        = {Pitfalls of Conformal Predictions for Medical Image Classification},
  author       = {Hendrik Alexander Mehrtens and Tabea-Clara Bucher and Titus Josef Brinker},
  year         = {2025},
  eprint       = {2506.18162},
  archiveprefix = {arXiv},
  primaryclass = {cs.LG},
  doi          = {10.48550/arXiv.2506.18162},
}

@misc{zhang2024truecam,
  title        = {Implementing Trust in Non-Small Cell Lung Cancer Diagnosis with a Conformalized Uncertainty-Aware {AI} Framework in Whole-Slide Images},
  author       = {Xiaoge Zhang and Tao Wang and Chao Yan and Fedaa Najdawi and Kai Zhou and Yuan Ma and Yiu-ming Cheung and Bradley A. Malin},
  year         = {2024},
  eprint       = {2501.00053},
  archiveprefix = {arXiv},
  primaryclass = {eess.IV},
  doi          = {10.1101/2024.12.27.24319715},
  note         = {medRxiv preprint},
}

@misc{fayjie2024predictive,
  title        = {Predictive uncertainty estimation in deep learning for lung carcinoma classification in digital pathology under real dataset shifts},
  author       = {Fayjie, Abdur R. and Borah, Jutika and Carbone, Florencia and Tack, Jan and Vandewalle, Patrick},
  year         = {2024},
  eprint       = {2408.08432},
  archiveprefix = {arXiv},
  primaryclass = {eess.IV},
  doi          = {10.48550/arXiv.2408.08432},
}

@article{hendrickx2024survey,
  author    = {Kilian Hendrickx and Lorenzo Perini and Dries Van der Plas and Wannes Meert and Jesse Davis},
  title     = {Machine Learning with a Reject Option: A Survey},
  journal   = {Machine Learning},
  volume    = {113},
  number    = {5},
  pages     = {3073--3110},
  year      = {2024},
  doi       = {10.1007/s10994-024-06534-x},
}
